% !PRINTED 51
\weberpage{51}
%<<CONT

\begin{equation}
 (x+y+z+\cdots)^n=\Pi(n)\sum^{\alpha,\nu}{x^\alpha u^\nu
\over \Pi(\alpha)\Pi(\nu)},
\tag{6}
\end{equation}
mit der Beschränkung
\begin{equation}
 \alpha+\nu=n.
\tag{7}
\end{equation}
Nun ist aber nach der Annahme schon bewiesen:
\begin{equation}
 u^\nu=\Pi(\nu)\sum^{\beta,\gamma,\ldots}
 {y^\beta z^\gamma\cdots
\over \Pi(\beta)\Pi(\gamma)\cdots},
\tag{8}
\end{equation}
\begin{equation}
 \beta+\gamma+\cdots=\nu,
\tag{9}
\end{equation}
und wenn dies in (6) eingesetzt wird, so ergiebt sich unmittelbar die Formel (3), und (7) geht in (4) über.

Die Coëfficienten
\begin{equation}
 P_{\alpha,\beta,\gamma,\ldots}^{(n)}={\Pi(n)\over
 \Pi(\alpha)\Pi(\beta)\Pi(\gamma)\cdots},
\tag{10}
\end{equation}
die ihrer Bedeutung nach ganze Zahlen sind, heissen die \gesperrt{Polynomialcoefficienten}.

Beispielsweise erhält man für die dritte Potenz des Trinoms:
\begin{equation}
\begin{aligned}
(x+y+z)^3={}&x^3+y^3+z^3+3x^2y+3xy^2+3x^2z+3xz^2\\
&+3y^2z+3yz^2+6xyz.
\end{aligned}
\tag{11}
\end{equation}

\begin{center}
{\large \S.\ 13.}

\vspace{0.5em}
{\large\gesperrt{Derivirte Functionen.}}
\end{center}

Es sei
\begin{equation}
 f(x)=a_0x^n+a_1x^{n-1}+a_2x^{n-2}+\cdots+a_n
\tag{1}
\end{equation}
eine ganze rationale Function $n$ter Ordnung.

Wenn wir darin $x$ durch ein Binom $x+y$ ersetzen, so können wir auf jedes einzelne Glied den binomischen Lehrsatz anwenden, und können das Ergebniss nach fallenden oder nach steigenden Potenzen von $x$ oder von $y$ ordnen. Wir wollen die Ordnung nach steigenden Potenzen von $y$ ausführen. Die höchste Potenz von $y$, die vorkommt, ist die $n$te, und der Coëfficient der nullten Potenz von $y$ ist die Function $f(x)$ selbst, wie man erkennt, wenn man $y=0$ setzt. Wir setzen also, indem wir die anderen Coëfficienten mit
\[
 f'(x),\qquad {f''(x)\over \Pi(2)},\qquad {f'''(x)\over \Pi(3)},\ldots
\]
bezeichnen:
%CONT>>
